\section{Choice of the construction constants}\label{sec:ordered-constants}

The estimates impose conditions on fixed constants and on the
running charge. The constants are chosen before the charge, mesh,
volume and twist. We first choose the polymer decay and the two
geometric scales; the charge bounds are then imposed for every
smaller charge. The local inverse constants are fixed before $M$,
whereas the weakened kernel constants may depend on $M$.
Neither group depends on the running charge. We verify this order
before imposing the small-charge conditions. The final subsection
uses the resulting construction to make the additional probe-dependent
restriction needed for the fourth cumulant.

\begin{proposition}\label{prop:ordered-constants}
Fix $0<\rho<1$ and the standard chamber weights.
The constants of the photon and fermion primitives,
small-field step, cover comparison and tuned flow can
be chosen in the order described below, with positive
charge ceilings $\bar e_0,\bar e_1$ independent of the mesh,
volume, twist and reference mass.
For every $0<\bar m\le1$, the physical ceiling $e_G$
in \eqref{eq:constants-physical-charge} is positive.
The terminal conditions \exteq{43}{TerminalPressure}
hold on the top and reference tori throughout
$0<e\le e_G$.
\end{proposition}

\subsection{The geometric scales and the positive exponent}

\begin{proof} Set $\bar\varepsilon=1/200$, $c_X=54$,
$\kappa_0=54\log72$, $K_0=72^{54}/36$,
$c_{\mathrm{rb}}=54K_0$ and $c_{\mathrm{fl}}=61c_XK_0$.
Choose
\[
 1-\bar\varepsilon<\alpha<1,\qquad
 \kappa\ge7\kappa_0+7.
\]
An odd integer $L\ge5$ is required to satisfy
\begin{align}
 L^{1/8-3\bar\varepsilon/2}&\ge2400,\qquad
 c_{\mathrm{fl}}L^{-3/8+6\bar\varepsilon}\le1/16,\notag\\
 L(\kappa-6\kappa_0-4)-1&\ge\kappa,\notag\\
 c_{\mathrm{rb}}L^3
       \exp\bigl(-L(\kappa-\kappa_0-1)+\kappa\bigr)
                                                   &\le1/2.
                                      \label{eq:constants-L}
\end{align}
These are simultaneous lower-bound conditions on $L$.
The powers $1/8-3\bar\varepsilon/2$ and
$3/8-6\bar\varepsilon$ are positive, as are
$\kappa-6\kappa_0-4$ and $\kappa-\kappa_0-1$.
The last expression tends to zero as $L\to\infty$.
Thus every sufficiently large odd $L$ satisfies
\eqref{eq:constants-L}.

After $L$ is fixed, the primitive construction fixes
$a'$, $C_V$, $\gamma_w$ and $\gamma_P$, with $a',C_V>0$
and $\gamma_w,\gamma_P>0$. Put
$N_c=(2L+1)^3$ and
$\widehat\kappa=N_c\kappa+\log(6c_X)>0$.
The photon and fermion
constants have no dependence on $M$ at this stage.
Put $\kappa_w=\kappa+\kappa_0+3$, $\alpha_0=1/54$ and $\kappa'=\kappa_w+\kappa_0$ and
\[
 M_F=\max\{4L+1,\ (108\pi C_Va')^2,\
                         108\pi C_Va'\exp(27\kappa')\}.
\]
Choose $M$ to be a power of $L$ satisfying
\begin{equation}\label{eq:constants-M}
 \begin{split}
 M&\ge L^2,\qquad M>4L,\qquad \gamma_wM\ge16,\\
 M^{\alpha_0}&\ge1+\exp(\widehat\kappa),\qquad
 M\ge M_F,\qquad 804\alpha_0\log M\le\gamma_PM.
 \end{split}
\end{equation}
All these conditions hold for every sufficiently large
power of $L$. For the last one, $\log x\le\sqrt x$
on $x\ge1$ gives the sufficient lower bound
$M\ge(804\alpha_0/\gamma_P)^2$.
The definition of $M_F$ gives the two fermion
conditions in \exteq{46}{FermionBlock},
because $27\alpha_0=1/2$.
The last inequality gives \exteq{38}{MaxwellBlock}.
The first four give the geometric separation, decay
and parameter-radius conditions used in localization.

The exponents can be fixed independently of the charge.
Take $r=2$, $p_0=9$, $p=19$ and
$1-\alpha<\varepsilon\le\bar\varepsilon$, and take $s=1/2$.
Then $p_0>3r+2$ and $p>2p_0$, while
\[
 \mathfrak e=s(1/4-9\varepsilon)-12\varepsilon
                      \ge17/400.
\]
For $N_*=8L^3$ and $a_L=(L^3-1)/(L-1)^3$,
\[
 \frac{2a_L}{N_*}
     <\frac1{4(L-1)^3}\le\frac1{256}<\frac{17}{400}.
\]
It follows that $\sigma=\mathfrak e-2a_L/N_*>0$.
In particular the positive exponent is fixed before
the charge is chosen.

\subsection{Calibration of the primitive and terminal bounds}

The chamber coercivity, heavy-sector gap and local inverse
construction fix $a'$ and $\gamma^F>0$ before $M$ is chosen;
see \exteq{34}{FermionBlock} and
\exteq{43}{FermionBlock}. The photon construction
likewise fixes $C_P^{\mathrm{ph}}$ and $\gamma_P>0$ before $M$;
its weakened families are considered only after
\exteq{38}{MaxwellBlock} holds.
The sums $c(\gamma)=\sum_{z\in\Z^3}
\exp(-\gamma|z|_\infty)$ are finite for $\gamma>0$.
Thus the constants used to choose $M$ in
\eqref{eq:constants-M} have already been fixed.

Keep $b=1$, $b_-=1-L^{-1}$ and $b_+=1$ from the inverse
construction; these are already fixed before $M$.
The coarse resolvent and its auxiliary block matrix give
\[
 C_{\mathrm{By}}=(1+2\exp(1))\sqrt2\max(1,b_-^{-1}),
 \qquad
 C_\Gamma=\frac{2(1+2\exp(1))}{b_-}
                  +128\exp(1)(1+2\exp(1))^2a'.
\]
The latter is the constant in
\exteq{62}{FermionBlock}. Define
\begin{equation}\label{eq:constants-fermion-calibration}
 \begin{split}
 C_1=\max\{&
 304\exp(1)\max(1,b)a'M^{27\alpha_0},
 C_\Gamma M^{27\alpha_0},\\
 &32\exp(1)\max(1,b^2)a'\exp(27\kappa')\},\\
 C^F=2\max\{&
 C_1,C_{\mathrm{By}},304\exp(1)\max(1,b)a',
 C_\Gamma,\exp(1)b^2(32a')^2\}.
 \end{split}
\end{equation}
The factor two converts Euclidean spinor bounds to
component bounds. The walk and coarse-family estimates
\exteq{47}{FermionBlock}--\exteq{48}{FermionBlock}
and \exteq{61}{FermionBlock}--\exteq{62}{FermionBlock}
are bounded by $C^F$ in this convention. The last entry
also bounds the compressed quadratic kernel.
At level zero the normalized kernel has the same bound,
and $M^{8\alpha_0}\le M^{27\alpha_0}$ bounds its diagonal
field multiplier.
The unweakened all-pair field bound
\exteq{60}{TerminalPressure} is at most
$608\exp(1)\max(1,b)a'\le C^F$; at zero its constant is
at most four, also below $C^F$.

Take the common primitive constants
\[
 C^P=116C_P^{\mathrm{ph}}M^{402\alpha_0},\qquad
 C_w^{(1)}=\max\{C^F,C^P,4\exp(2\gamma_w)M^{27\alpha_0}\},
 \qquad C_w=64K_0C_w^{(1)}.
\]
These are precisely \exteq{23}{CoverTuning}.
The top systems and their projections therefore satisfy
the same bounds at every required level and side.
The determinant, photon-domain and localization constants
are given in Sections~\extref{6.2}{FermionBlock},
\extref{6}{MaxwellBlock} and
\extref{5.1}{NonlinearStep}. In particular,
\[
 \begin{aligned}
 c_C&=3L^{9/2},\qquad C_W=28C_w^2c_\gamma,\\
 C_d=16c_XC_w
       &\max\{\tfrac\pi2b\exp(1)(1+2\exp(1)),
                                      2+58\exp(1)+b\exp(1)\}.
 \end{aligned}
\]
The kernel-sum constant $c_\gamma$ is bounded by
$2L^3\exp((2L+1)\gamma_w)c(\gamma_w)$;
the local-gauge constant $C_B$ depends only on
$C_P^{\mathrm{ph}},\gamma_P,L,M$.
All these constants are finite after $M$ is fixed.

For the terminal norms fix
\begin{equation}\label{eq:constants-terminal-calibration}
 C_0=\max\{C_P^{\mathrm{ph}}c_{\mathrm{bd}}(\gamma_P),
                 C^F,4b_++16C^Fc(\gamma^F)\},
 \qquad c_{\mathrm{bd}}(\gamma)=3\exp(3\gamma)c(\gamma).
\end{equation}
The photon row sums, both field-map row norms and the
terminal Dirac form in
Section~\extref{6.2}{TerminalPressure} are then
bounded by this $C_0$. In particular $C_0\ge1$.
The polymer sum has $1\le c_\kappa\le K_0$.
The free Bloch construction fixes the twist constants
$C_\Sigma<\infty$ and $\gamma_\Sigma>0$ independently of
the running charge, mass in its prescribed range, mesh
and side; see \exteq{37}{TerminalPressure}.

The differentiated-response constants are now fixed by
the terminal inverse and field-map estimates:
$\mathcal S=C^F(24/\gamma^F)^3$,
$C_K=226\mathcal S^2$, $C_K'=1328\mathcal S^2$ and
$C_H=203M(C^F)^2$. The real covariance coefficients
$g,h_G,g_1,\alpha_G$ and its rate $\vartheta_0$ are
defined with \exteq{74}{TerminalPressure}.
The insertion and mixed-logarithm coefficients are
$C_V^{\mathrm{ins}}=10^3C_H^2(24/\gamma^F)^{12}K_0$ and
$C_W^{\mathrm{mix}}$ in \exteq{109}{ReferenceResponse}.
Their inverse-decay rates are positive and their
coefficients finite because $b_->0$,
$\mu_-=L^{-3}>0$ and $\gamma^F>0$.
They depend on $L,M$ and the fixed chamber data;
they have no dependence on the later charge or probe.
The probe-dependent constants will be chosen last.

\subsection{Uniform running-charge bounds}

All coefficients of the small-charge conditions are
fixed after the preceding choices. The fluctuation substitution
constant from \cite[Section~1]{NonlinearStep} is
\[
 C_T=2\exp(3/2)
       [1+b_-^{-1}(54L\exp(6)+12\exp(1)+1)].
\]
We write $\bar e_0,\bar e_1$ for the fixed charge ceilings,
to distinguish them from the level-zero running charge $e_0$.
They are the ceilings denoted $e_0,e_1$ in the terminal and
response estimates, and agree with the barred convention in
\cite{CoverTuning}. The elementary
estimate used to meet them simultaneously is as follows.
For finitely many $A_i,d_i,\lambda_i$ with
$A_i,d_i\ge0$ and $\lambda_i>0$, and for $b>0$, put
\begin{equation}\label{eq:constants-tail-threshold}
 V(A,d,\lambda;b)
   =\max_i\frac2{\lambda_i}
       \log^+\left[
         \frac{nA_i}{b}
         \left(\frac{2d_i}{\lambda_i\exp(1)}\right)^{d_i}
       \right].
\end{equation}
Here $n$ is the number of terms, $\log^+0=0$
and a factor with $d_i=0$ is interpreted as one.
For all $v\ge V(A,d,\lambda;b)$ with $v>0$,
\begin{equation}\label{eq:constants-tail-bound}
 \sum_i A_iv^{d_i}\exp(-\lambda_iv)\le b.
\end{equation}
Indeed $v^d\exp(-\lambda v/2)$ has maximum
$(2d/(\lambda\exp(1)))^d$ for $d>0$.
For $d=0$ use $\exp(-\lambda v)\le\exp(-\lambda v/2)$.
The threshold makes each resulting term at most $b/n$.

The primitive domains, determinant, adjustment and nonlinear
step require \exteq{43}{MaxwellBlock},
\exteq{8}{NormalizedContraction},
\exteq{21}{NormalizedContraction},
\exteq{1}{NormalizedContraction},
\exteq{81}{FermionBlock} and
\exteq{3}{NonlinearStep}.
Their charge conditions are obtained from the following
finite families for $v=-\log e_k$.
Write $T(A,d,\lambda)=Av^d\exp(-\lambda v)$.
The polynomial-exponential conditions are
\begin{align}
 T(1,p_0,\varepsilon/4)&\le1,&
 T(1,0,5\varepsilon/2)&\le\tfrac12,\notag\\
 T(c_{\mathrm{fl}},0,\varepsilon)&\le1,&
 T(3M,0,\varepsilon)&\le1,\notag\\
 T(2C_w^2c_\gamma^2c_C,0,\varepsilon/4)&\le1,&
 T(16c_Xc_\gamma C_wC_T^2M^3,0,\varepsilon/2)&\le1,\notag\\
 T(96M^3,0,1/4)&\le1,&
 T(32C_wC_T,0,3\varepsilon/2)&\le1,\notag\\
 T(88M^3,0,1/4-6\varepsilon)&\le1,&
 T(112C_w^2c_\gamma,0,2\varepsilon)&\le3,\notag\\
 T(c_S,0,1/4-5\varepsilon)&\le1,&
 T(c_SL^3,0,\varepsilon)&\le1,\notag\\
 T(36c_X^2K_0^2,0,1/4-5\varepsilon)&\le1,&
 T(1,0,\varepsilon)&\le\tfrac14,\notag\\
 T(3M,0,1/4+\varepsilon)+T(3M,0,2\varepsilon)&\le1,\notag\\
 T(1,0,1/4)+T(1,0,\varepsilon)&\le1,\notag\\
 T(160L^{1+3\varepsilon/2}M,0,\varepsilon)&\le1,&
 T(90c_XLM,0,\varepsilon)&\le1,\notag\\
 T(9ML^\varepsilon,0,2\varepsilon)&\le1,\notag\\
 T(2M,1,\varepsilon)+T(32M,0,\varepsilon)&\le1,\notag\\
 T(8C_B,p,3/4-3\varepsilon)&\le1,&
 T(8,0,3/4-7\varepsilon/2)&\le1,\notag\\
 T(6L,0,\varepsilon)&\le1,&
 T(4C_Va',0,\varepsilon)&\le1,\notag\\
 T(4,0,1/4+2\varepsilon)+T(2,0,3\varepsilon)&\le1,\notag\\
 T(8c_{\mathrm{rb}}L^3C_dM^3,0,2\varepsilon)&\le1.
                                      \label{eq:constants-running-inventory}
\end{align}
Here $c_S=10^5c_X^5K_0^5\exp(6\kappa)$.
The inequality with $96M^3$ controls the first step;
the one with $C_T^2M^3$ controls the positive-level
fluctuation substitution. The conditions with $c_S$
and $36c_X^2K_0^2$ control the connected logarithm and its scaling.
The domain conditions occur with target $1/4$,
with rates $(1/4,\varepsilon)$, with coefficients
$8C_B,8$, and with coefficients $4,2$.
The remaining displayed charge hypotheses above are
individual inequalities of this inventory;
\exteq{83}{NonlinearStep} follows from the
condition containing $c_S$.

Every decay rate in the inventory is positive for
$0<\varepsilon\le1/200$. The only additional
running-charge restrictions in this group are
\[
 v^{2p_0}/2-7v\ge\log12,\qquad v^2\ge10L,\qquad v\ge9.
\]
They hold respectively for $v\ge15$,
$v\ge\sqrt{10L}$ and $v\ge9$.
Taking the maximum of these thresholds and the finitely
many values in \eqref{eq:constants-tail-threshold}
gives a finite $V_0$. Choose a running-charge ceiling
$0<\bar e_0<\exp(-V_0)$ so that every such condition
holds for every $v\ge-\log \bar e_0$, equivalently for
every running charge in $(0,\bar e_0]$.
The set of permitted ceilings is downward closed.

The terminal-scale conditions require a second ceiling
$\bar e_1\le\min(\bar e_0,1/3)$. Write $\ell(v)=LMv^r$.
They follow from
\begin{align}
 (c_{\mathrm{fl}}+2)\exp(-\varepsilon v)
                      &\le\tfrac1{16}(1-L^{-\varepsilon/2}),\notag\\
 8C_0^2\ell(v)^6(b_-^{-1}+L\ell(v))
                      &\le\exp(\varepsilon v/2),\notag\\
 4(2C_0+c_\kappa+5c_{\mathrm{fl}}+6)\ell(v)^3
                      &\le\exp(\varepsilon v),\notag\\
 \exp(-(1/4-9\varepsilon)v)&\le\tfrac18,\notag\\
 4C_0\ell(v)^2v^p&<\exp(\varepsilon v),\notag\\
 12\ell(v)^3\left(\exp(-v^{p_0}/2)
                    +\exp(-v^{2p}/(2C_0^2))\right)
                      &\le\exp(-v),\notag\\
 C_\Sigma\exp(-\gamma_\Sigma Mv^r)&\le\exp(-v),\notag\\
 L^2\exp(-2v)\ell(v)&\le1,\notag\\
 6N_*Mv^r\exp(-v)(v\exp(-v)+\tfrac12)&\le\pi^2.
                                      \label{eq:constants-terminal-inventory}
\end{align}
These are the scalar hypotheses of
\exteq{48}{CoverTuning},
\exteq{43}{TerminalPressure} and
\eqref{eq:transfer-tail-condition}, after substituting
\exteq{51}{CoverTuning}.
Every polynomial-exponential rate here is positive for
$0<\varepsilon\le1/200$. The reference flux condition expands as
\[
          6N_*Mv^{r+1}\exp(-2v)+3N_*Mv^r\exp(-v)\le\pi^2.
\]
The superlinear tails reduce to these bounds:
for $v\ge\max(4,4C_0^2)$,
\[
 v-\tfrac12v^{p_0}\le-v,\qquad
 v-\frac{v^{2p}}{2C_0^2}\le-v.
\]
Hence
\[
 12\ell(v)^3
 \left[\exp\left(v-\tfrac12v^{p_0}\right)
       +\exp\left(v-\frac{v^{2p}}{2C_0^2}\right)\right]
                \le24(LM)^3v^{3r}\exp(-v).
\]
For the twist tail, if
$v\ge\max(1,2/(\gamma_\Sigma M),\log^+C_\Sigma)$,
then $C_\Sigma\exp(v-\gamma_\Sigma Mv^r)\le1$.
Put $\mathrm{lg}=\tfrac12\log L$.
The scale-ratio condition follows for $v\ge2r$:
\begin{equation}\label{eq:constants-scale-ratio}
 \left(1+\frac{3\mathrm{lg}}{v}\right)^r
        \le\exp\left(\frac{3r\mathrm{lg}}{v}\right)
        \le L^{3/4}<L.
\end{equation}
Together with $v\ge8$, $v\ge31L+17$ and
$p_0v^{p_0-1}\ge28$, these bounds give a finite
threshold $V_1$. Its maximum includes $2r$.
Let $\mathcal E_1$ consist of the numbers
$e'\in(0,\min(\bar e_0,1/3)]$ for which the geometric
conditions and \eqref{eq:constants-terminal-inventory}
hold for every $v\ge-\log e'$. It contains
$\min(\bar e_0,1/3,\exp(-V_1))>0$ and is downward closed.
Fix
\[
 \bar e_1=\tfrac12\sup\mathcal E_1>0.
\]
There is an element of $\mathcal E_1$ larger than $\bar e_1$;
downward closure therefore gives $\bar e_1\in\mathcal E_1$.
Thus all terminal-scale conditions hold for every
$v\ge-\log \bar e_1$.

\subsection{The physical charge interval}

Fix the reference mass and hence $h\ge1$. Set
\begin{equation}\label{eq:constants-physical-charge}
 e_G=
 \begin{cases}
 \bar e_1,&h=1,\\
 \min\{\bar e_1/\sqrt h,\,1/(\sqrt h\log h)\},&h>1.
 \end{cases}
\end{equation}
This is positive. For $0<e\le e_G$, every running charge
up to level $K$ is at most $\bar e_1$, and
$e\sqrt h\log h\le1$.
The scale proposition therefore applies on the top and
reference tori. The terminal inventory gives
\exteq{43}{TerminalPressure} at their respective stops,
and \eqref{eq:transfer-W3} holds on the reference torus.
The common pair of Theorem~\extref{D}{CoverTuning}
is defined for $N\ge N_{\min}(R_{\mathrm{top}},e)$.
At lower cutoffs only its mass shift is assigned zero.
For $N\ge N_{\mathrm{ref}}(e)$ the same pair is
controlled on the reference torus. Its mass shift is real,
independent of the cutoff volume, and has absolute value
less than $1/8$.
\end{proof}

\subsection{The probe-dependent charge interval}

\begin{proposition}\label{prop:ordered-probe-constants}
Fix the constants of \cref{prop:ordered-constants} and a compact
real probe $\beta$ with $b(dd^*\beta)>0$. Let $C_U$ be the constant
in \exteq{59}{VolumeTransfer}, depending only on the fixed construction,
reference mass and probe, and set $C_f=8C_U+3C_U^2$.
There is a positive $e_\beta\le e_G$ such that, for every
$0<e\le e_\beta$ and $N\ge N_{\rm ref}(e)$,
\eqref{eq:transfer-small-Wick} holds with $\sigma>0$ and
$C_fe^\sigma\le b(dd^*\beta)/2$.
\end{proposition}

\begin{proof} Fix a compact real probe $\beta$. Its buffered support
diameter, energy and uniform rescaled norms are finite by
\exteq{59}{ReferenceResponse}.
The differentiated inverse and interaction constants have
already been fixed; these probe constants are chosen last.
Proposition~\extref{11.1}{ReferenceResponse}
gives $e_\sharp(\beta)>0$ for which all its factor,
one-cell, tree, jet and placement conditions hold for
every $0<e\le e_\sharp(\beta)$, uniformly in
$N\ge N_{\mathrm{ref}}(e)$ and in the allowed placements.
Its proof uses $a_*\in\{h/L,h/L^2\}$ and
$v_*=v_K+O(1)$: polynomial factors have fixed degrees,
the separated-cell errors decay exponentially in $v_*^r$,
and the complement decays faster than the Cauchy-radius
loss. This gives one interval of charges for all cutoffs.

For $0<e\le\min(e_G,e_\sharp(\beta))$,
\exteq{5}{TerminalPressure} and
\exteq{119}{ReferenceResponse} supply the
two inputs in \eqref{eq:transfer-terminal-input}.
The probe placement inequalities imply
\exteq{12}{ReferenceResponse} for $\eta\le1$.
The bounds for the free coefficients and $C_U$ in
\exteq{59}{VolumeTransfer} are finite and independent of
$e,N,T$ and placement. They depend only on the fixed
construction and probe.

Suppose now that $b(dd^*\beta)>0$ and put
$C_f=8C_U+3C_U^2$.
Choose
\begin{equation}\label{eq:constants-probe-charge}
 0<e_\beta\le
 \min\{e_G,e_\sharp(\beta),1,
                         (b(dd^*\beta)/(2C_f))^{1/\sigma}\}.
\end{equation}
The minimum is positive, since $\sigma>0$ and $C_f>0$.
\Cref{prop:transfer-UV-inputs} holds throughout
$(0,e_\beta]$ and gives the finite-cutoff hypotheses of
Theorem~\extref{A}{VolumeTransfer}. Its component clause will be
applied after the component Gaussian bound has been established.
In particular $C_fe^\sigma\le b(dd^*\beta)/2$.
The construction constants were chosen before $h$ and
the probe; the physical charge interval precedes the
probe-dependent restriction.
\end{proof}


